\documentclass[11pt,a4paper]{article}
\usepackage[margin=2.2cm]{geometry}
\usepackage{amsmath,amssymb}
\usepackage{enumitem}
\usepackage{fancyhdr}
\usepackage{booktabs}
\usepackage{array}
\usepackage{pifont}
\usepackage[most]{tcolorbox}
\usepackage{kotex}
\setmainhangulfont{Nanum Myeongjo}
\setsanshangulfont{Apple SD Gothic Neo}
\usepackage[hidelinks]{hyperref}

\pagestyle{fancy}
\fancyhf{}
\lhead{Algorithms $\times$ pandas}
\rhead{알고리즘으로 보는 pandas와 Datathon}
\cfoot{\thepage}
\setlength{\headheight}{14pt}
\setlength{\parindent}{0pt}
\setlength{\parskip}{0.45em}
\setlength{\emergencystretch}{2em}
\setlist[itemize]{leftmargin=1.3em,itemsep=0.2em,topsep=0.2em}
\setlist[enumerate]{leftmargin=1.6em,itemsep=0.2em,topsep=0.2em}
\renewcommand{\arraystretch}{1.15}
\AddToHook{cmd/section/before}{\clearpage}
\sloppy

\newcommand{\term}[1]{\textbf{#1}}
\newcommand{\code}[1]{\texttt{#1}}
\newcommand{\aref}[1]{{\footnotesize\textsf{\color{blue!55!black}[Alg #1]}}}
\newcommand{\sref}[1]{{\footnotesize\textsf{\color{green!40!black}[통계책 #1]}}}
\newcommand{\aside}[1]{{\footnotesize\color{black!65}#1}}
\newcommand{\Ans}{\par\textbf{Answer.}\ }
\newcommand{\Just}{\par\textbf{Justification.}\ }

\newtcolorbox{codebox}{breakable,colback=black!3,colframe=black!45,boxrule=0.5pt,arc=1.5pt,
  left=4pt,right=4pt,top=2pt,bottom=2pt,fontupper=\small\ttfamily}
\newtcolorbox{keybox}[1]{breakable,colback=orange!4,colframe=orange!60!black,boxrule=0.6pt,arc=1.5pt,
  fonttitle=\bfseries\sffamily,title={#1}}
\newtcolorbox{dtbox}{breakable,colback=blue!3,colframe=blue!45!black,boxrule=0.6pt,arc=1.5pt,
  fonttitle=\bfseries\sffamily,title={Datathon에서는}}

\begin{document}

\begin{center}
{\LARGE\textbf{알고리즘으로 보는 pandas와 Datathon}}\\[0.4em]
{\large 정렬 · 힙 · 탐색 · 그래프 알고리즘은 데이터 분석의 어디에 숨어 있나}\\[0.3em]
{\small \aref{13 Heap sort} = Algorithms 강의 슬라이드 번호, \sref{3장} = ``pandas 흐름으로 따라가는 통계'' 문서의 장.\\
시간은 이 컴퓨터(Apple Silicon, numpy 2.1, pandas 2.2)에서 3회 중 최솟값으로 직접 측정했습니다.}
\end{center}

\vspace{0.6em}
\begin{keybox}{결론부터: 쓸 일이 있나?}
\textbf{직접 구현할 일은 거의 없지만, pandas·numpy·sklearn 내부에서 이미 쓰이고 있습니다.} 알고 있으면 (1) 왜 어떤 코드가 수백 배 느린지 이해하고, (2) 함수의 옵션(\code{kind=}, \code{max\_leaf\_nodes=})을 의미 있게 고르고, (3) 결과(트리 모형, 분위수)를 정확히 해석할 수 있습니다.
\begin{center}\small
\begin{tabular}{@{}l l l c@{}}
\toprule
배운 것 & pandas/라이브러리에서 쓰이는 곳 & 장 & 관련도 \\
\midrule
Big-O, RAM 모델 & 반복문 vs 벡터화, \code{iterrows} 금지 & 1 & ★★★ \\
Quick sort + Heap sort & \code{sort\_values} 기본값 = introsort (둘의 조합) & 2 & ★★★ \\
Merge sort, 안정성 & \code{kind="stable"} (timsort), 다중 키 정렬 & 2 & ★★★ \\
Counting / Radix sort & \code{groupby} 내부의 counting sort, 정수 정렬 & 2, 4 & ★★☆ \\
정렬 하한 $\Omega(n\log n)$ & 왜 radix가 더 빠를 수 있나 & 2 & ★☆☆ \\
Selection (quickselect) & \code{median}, \code{quantile}, \code{nlargest} & 3 & ★★★ \\
Heap & top-$k$, 결정트리 best-first, Dijkstra & 3, 6, 9 & ★★☆ \\
Binary search & \code{searchsorted}, \code{pd.cut}, \code{merge\_asof} & 5 & ★★★ \\
Hashing & \code{groupby}, \code{merge}, \code{value\_counts} & 4 & ★★★ \\
DFS / BFS, 트리 & 결정트리 학습·출력, 연결요소(중복 제거) & 6 & ★★☆ \\
Heuristic, greedy & 결정트리 분할, 변수 선택 (stepwise) & 7 & ★★☆ \\
Minimax & 직접은 없음. 개념: minimax 추정, minimax 공정성 & 8 & ★☆☆ \\
Dijkstra & 이번 datathon엔 없음. 네트워크·교통 데이터용 & 9 & ☆☆☆ \\
Divide \& conquer, 점화식 & split-apply-combine, bootstrap 비용 & 10 & ★☆☆ \\
\bottomrule
\end{tabular}
\end{center}
\aside{※ 4시간짜리 채용 bias datathon에서 실제로 \emph{판단에 영향을 주는} 것은 1–7장입니다. 8–10장은 ``어디에 쓰이는지 궁금할 때'' 읽으면 됩니다.}
\end{keybox}

\clearpage
\clearpage
% 세 책 공통: 앞쪽 "전체 흐름" 페이지. 각 책에서 % 세 책 공통: 앞쪽 "전체 흐름" 페이지. 각 책에서 % 세 책 공통: 앞쪽 "전체 흐름" 페이지. 각 책에서 \input{roadmap} 으로 넣는다.
\providecommand{\code}[1]{\texttt{#1}}
\newcolumntype{L}[1]{>{\raggedright\arraybackslash}p{#1}}
% 책별 장 표시 (작은 색 주석)
\newcommand{\Wc}[1]{{\scriptsize\textsf{\color{blue!60!black}[흐름 #1]}}}
\newcommand{\Hc}[1]{{\scriptsize\textsf{\color{orange!70!black}[응용 #1]}}}
\newcommand{\Ac}[1]{{\scriptsize\textsf{\color{green!45!black}[알고 #1]}}}
\newcommand{\Pc}[1]{{\scriptsize\textsf{\color{red!60!black}[실습 #1]}}}

\section*{먼저 읽기: 분석 전체 흐름과 책 안내}
\addcontentsline{toc}{section}{먼저 읽기: 분석 전체 흐름과 책 안내}

\textbf{네 권의 역할}
\begin{center}\small
\begin{tabular}{@{}l L{5.4cm} L{7.4cm}@{}}
\toprule
표시 & 책 (파일) & 역할 \\
\midrule
\Wc{} & pandas 흐름으로 따라가는 통계 (\code{pandas-stats-walkthrough}) & \textbf{기초}: 평균·SD·정규·상관·회귀·검정을 pandas 순서대로 \\
\Hc{} & 채용 데이터로 다시 배우는 통계 (\code{hiring-statistics-book}) & \textbf{응용}: 통제 회귀, logistic, Bayesian, 공정성으로 결론 내기 \\
\Ac{} & 알고리즘으로 보는 pandas (\code{algorithms-for-datathon}) & \textbf{도구}: 빠르고 정확한 코드, 결정트리·greedy 보조 분석 \\
\Pc{} & numpy·pandas로 직접 만드는 통계 (\code{stats-basics-workbook}) & \textbf{실습}: 빈칸 노트북으로 통계 공식을 직접 구현, 필요한 곳에 알고리즘 적용 \\
\bottomrule
\end{tabular}
\end{center}
\aside{※ 아래 표의 작은 색 주석이 ``이 단계는 어느 책 몇 장''입니다. 예: \Wc{5} = 흐름 책 5장, \Hc{6.9} = 응용 책 6.9절, \Ac{4} = 알고리즘 책 4장. 0장 = 응용 책 앞의 ``데이터 미리보기와 pandas 작동 원리''.}

\vspace{0.4em}
\textbf{Datathon 분석 흐름: 데이터를 받은 순간부터 발표까지}
\begin{center}\small
\renewcommand{\arraystretch}{1.25}
\begin{tabular}{@{}c L{3.3cm} L{4.6cm} L{5.6cm}@{}}
\toprule
단계 & 하는 일 & 대표 pandas 코드 & 해당 장 \\
\midrule
① & 질문을 통계 질문으로 & -- &
  \Hc{1} \Wc{1 추정·가설·검증 상자} \\
② & 불러오고 구조 확인 & \code{read\_csv}, \code{head}, \code{info} &
  \Hc{0.1–0.3} \Wc{1} \\
③ & 변수 분류 (결과 $Y$ / 보호속성 $A$ / 자격 $X$) & \code{dtypes}, \code{value\_counts} &
  \Hc{2} \Wc{1.1} \\
④ & 결측·이탈값 처리 & \code{isna}, \code{dropna}, \code{quantile} &
  \Hc{0.7, 2.4} \Wc{1.3, 2.3} \\
⑤ & 분포 요약 (중심, 퍼짐, 정규성) & \code{describe}, \code{std}, \code{(x-x.mean())/x.std()} &
  \Wc{2, 3} \Hc{3.1} \Ac{3 분위수} \Pc{1–3} \\
⑥ & 그룹별 합격률과 신뢰구간 & \code{groupby().mean()} &
  \Hc{0.5, 3.2} \Wc{4} \Ac{4 groupby 내부} \Pc{4} \\
⑦ & 교란 변수 확인 (Simpson) & \code{crosstab}, \code{groupby([a,b])} &
  \Hc{3.3} \Wc{12 ⑧} \\
⑧ & 가설검정 (차이가 우연인가) & \code{proportions\_ztest}, \code{chi2\_contingency}, \code{ttest\_ind} &
  \Wc{5, 11, 12} \Hc{5} \Pc{5–7} \\
⑨ & 상관과 단순회귀 & \code{corr}, \code{pd.cut}, \code{smf.ols} &
  \Wc{6–8} \Ac{5 pd.cut} \Pc{8–9} \\
⑩ & 통제 회귀와 interaction (핵심) & \code{smf.ols("y \~{} A + X")}, \code{A*B} &
  \Hc{0.8–0.9, 6} \Wc{13} \Pc{10} \\
⑪ & 다른 방법으로 재확인 & \code{smf.logit}, \code{rng.beta}, \code{DecisionTreeClassifier} &
  \Hc{7, 8} \Wc{9 bootstrap} \Pc{6, 11} \\
⑫ & 공정성 지표와 한계 & \code{rate["F"]/rate["M"]}, \code{idxmin} &
  \Hc{9} \Ac{8 minimax} \\
⑬ & 정리와 발표 & -- &
  \Hc{9.2 결론 문장, 11 체크리스트} \Wc{14 요약} \Ac{11} \\
\bottomrule
\end{tabular}
\end{center}
\aside{※ 이론이 막히면: 추정량·MoM·MLE \Wc{9–10} \Hc{4}, 가설 세우기·검정력 \Wc{11}, 코드가 느리면 \Ac{1–2, 4}. 한 번에 돌려 보기: \code{python3 algo\_toolkit.py hiring\_practice.csv} (⑥⑦⑨⑪⑫ 요약).}

\clearpage
\textbf{6시간 읽기 계획 (읽기 편한 순서)}
\begin{center}\small
\renewcommand{\arraystretch}{1.25}
\begin{tabular}{@{}l L{4.6cm} L{6.6cm} r@{}}
\toprule
시간 & 무엇을 & 어디를 & 분 \\
\midrule
0:00 & 데이터와 pandas 감 잡기 & \Hc{0} 전체, \Hc{1} 전체 지도 & 30 \\
0:30 & 기초 통계 1: 요약과 분포 & \Wc{1–3} (\code{describe}, 표준화, 정규근사) & 50 \\
1:20 & 기초 통계 2: 추정과 검정 & \Wc{4–5} (신뢰구간, $z$/$\chi^2$/$t$) & 40 \\
2:00 & \textit{휴식} & & 10 \\
2:10 & 기초 통계 3: 상관과 회귀 & \Wc{6–8} (회귀 효과, RMSE, 잔차도) & 50 \\
3:00 & 검정 고르는 법 & \Wc{12} 검정 도구상자 (표 + 관심 있는 카드만) & 30 \\
3:30 & 응용 1: 통제와 interaction & \Hc{3, 6} + 각 장 끝 ``pandas 코드'' 직접 실행 & 50 \\
4:20 & \textit{휴식} & & 10 \\
4:30 & 응용 2: 재확인 & \Hc{7} logistic, \Hc{8} Bayesian (8.4 ``왜''는 꼭) & 40 \\
5:10 & 결론과 한계 & \Hc{9}, \Hc{11} 체크리스트, \Ac{표지 결론표} & 20 \\
5:30 & 도구 실행과 발표 문장 연습 & \code{algo\_toolkit.py} 실행, \Hc{9.2} 문장을 내 말로 & 30 \\
(다음 날) & 코드로 손에 익히기 & \Pc{0–12} 빈칸 노트북 (\code{stats\_basics\_practice.ipynb}) & 120 \\
\bottomrule
\end{tabular}
\end{center}
\aside{※ 범위 밖 (학부 1–2학년 목표를 넘는 것): sklearn을 쓰는 \Ac{6–7}의 결정트리와 \code{advanced/} 폴더의 분류 워크북은 나중에 선택으로.}\\
\aside{※ 건너뛰어도 되는 곳 (나중에 참고용): \Wc{9–11, 13}, \Hc{4, 10}, \Ac{1–10} 본문, 모든 연습 문제. 시간이 남으면 \Wc{15} 연습 문제 1–8번부터.}\\
\aside{※ 각 장은 새 페이지에서 시작하고, 목차의 항목을 누르면 그 쪽으로 이동합니다.}
 으로 넣는다.
\providecommand{\code}[1]{\texttt{#1}}
\newcolumntype{L}[1]{>{\raggedright\arraybackslash}p{#1}}
% 책별 장 표시 (작은 색 주석)
\newcommand{\Wc}[1]{{\scriptsize\textsf{\color{blue!60!black}[흐름 #1]}}}
\newcommand{\Hc}[1]{{\scriptsize\textsf{\color{orange!70!black}[응용 #1]}}}
\newcommand{\Ac}[1]{{\scriptsize\textsf{\color{green!45!black}[알고 #1]}}}
\newcommand{\Pc}[1]{{\scriptsize\textsf{\color{red!60!black}[실습 #1]}}}

\section*{먼저 읽기: 분석 전체 흐름과 책 안내}
\addcontentsline{toc}{section}{먼저 읽기: 분석 전체 흐름과 책 안내}

\textbf{네 권의 역할}
\begin{center}\small
\begin{tabular}{@{}l L{5.4cm} L{7.4cm}@{}}
\toprule
표시 & 책 (파일) & 역할 \\
\midrule
\Wc{} & pandas 흐름으로 따라가는 통계 (\code{pandas-stats-walkthrough}) & \textbf{기초}: 평균·SD·정규·상관·회귀·검정을 pandas 순서대로 \\
\Hc{} & 채용 데이터로 다시 배우는 통계 (\code{hiring-statistics-book}) & \textbf{응용}: 통제 회귀, logistic, Bayesian, 공정성으로 결론 내기 \\
\Ac{} & 알고리즘으로 보는 pandas (\code{algorithms-for-datathon}) & \textbf{도구}: 빠르고 정확한 코드, 결정트리·greedy 보조 분석 \\
\Pc{} & numpy·pandas로 직접 만드는 통계 (\code{stats-basics-workbook}) & \textbf{실습}: 빈칸 노트북으로 통계 공식을 직접 구현, 필요한 곳에 알고리즘 적용 \\
\bottomrule
\end{tabular}
\end{center}
\aside{※ 아래 표의 작은 색 주석이 ``이 단계는 어느 책 몇 장''입니다. 예: \Wc{5} = 흐름 책 5장, \Hc{6.9} = 응용 책 6.9절, \Ac{4} = 알고리즘 책 4장. 0장 = 응용 책 앞의 ``데이터 미리보기와 pandas 작동 원리''.}

\vspace{0.4em}
\textbf{Datathon 분석 흐름: 데이터를 받은 순간부터 발표까지}
\begin{center}\small
\renewcommand{\arraystretch}{1.25}
\begin{tabular}{@{}c L{3.3cm} L{4.6cm} L{5.6cm}@{}}
\toprule
단계 & 하는 일 & 대표 pandas 코드 & 해당 장 \\
\midrule
① & 질문을 통계 질문으로 & -- &
  \Hc{1} \Wc{1 추정·가설·검증 상자} \\
② & 불러오고 구조 확인 & \code{read\_csv}, \code{head}, \code{info} &
  \Hc{0.1–0.3} \Wc{1} \\
③ & 변수 분류 (결과 $Y$ / 보호속성 $A$ / 자격 $X$) & \code{dtypes}, \code{value\_counts} &
  \Hc{2} \Wc{1.1} \\
④ & 결측·이탈값 처리 & \code{isna}, \code{dropna}, \code{quantile} &
  \Hc{0.7, 2.4} \Wc{1.3, 2.3} \\
⑤ & 분포 요약 (중심, 퍼짐, 정규성) & \code{describe}, \code{std}, \code{(x-x.mean())/x.std()} &
  \Wc{2, 3} \Hc{3.1} \Ac{3 분위수} \Pc{1–3} \\
⑥ & 그룹별 합격률과 신뢰구간 & \code{groupby().mean()} &
  \Hc{0.5, 3.2} \Wc{4} \Ac{4 groupby 내부} \Pc{4} \\
⑦ & 교란 변수 확인 (Simpson) & \code{crosstab}, \code{groupby([a,b])} &
  \Hc{3.3} \Wc{12 ⑧} \\
⑧ & 가설검정 (차이가 우연인가) & \code{proportions\_ztest}, \code{chi2\_contingency}, \code{ttest\_ind} &
  \Wc{5, 11, 12} \Hc{5} \Pc{5–7} \\
⑨ & 상관과 단순회귀 & \code{corr}, \code{pd.cut}, \code{smf.ols} &
  \Wc{6–8} \Ac{5 pd.cut} \Pc{8–9} \\
⑩ & 통제 회귀와 interaction (핵심) & \code{smf.ols("y \~{} A + X")}, \code{A*B} &
  \Hc{0.8–0.9, 6} \Wc{13} \Pc{10} \\
⑪ & 다른 방법으로 재확인 & \code{smf.logit}, \code{rng.beta}, \code{DecisionTreeClassifier} &
  \Hc{7, 8} \Wc{9 bootstrap} \Pc{6, 11} \\
⑫ & 공정성 지표와 한계 & \code{rate["F"]/rate["M"]}, \code{idxmin} &
  \Hc{9} \Ac{8 minimax} \\
⑬ & 정리와 발표 & -- &
  \Hc{9.2 결론 문장, 11 체크리스트} \Wc{14 요약} \Ac{11} \\
\bottomrule
\end{tabular}
\end{center}
\aside{※ 이론이 막히면: 추정량·MoM·MLE \Wc{9–10} \Hc{4}, 가설 세우기·검정력 \Wc{11}, 코드가 느리면 \Ac{1–2, 4}. 한 번에 돌려 보기: \code{python3 algo\_toolkit.py hiring\_practice.csv} (⑥⑦⑨⑪⑫ 요약).}

\clearpage
\textbf{6시간 읽기 계획 (읽기 편한 순서)}
\begin{center}\small
\renewcommand{\arraystretch}{1.25}
\begin{tabular}{@{}l L{4.6cm} L{6.6cm} r@{}}
\toprule
시간 & 무엇을 & 어디를 & 분 \\
\midrule
0:00 & 데이터와 pandas 감 잡기 & \Hc{0} 전체, \Hc{1} 전체 지도 & 30 \\
0:30 & 기초 통계 1: 요약과 분포 & \Wc{1–3} (\code{describe}, 표준화, 정규근사) & 50 \\
1:20 & 기초 통계 2: 추정과 검정 & \Wc{4–5} (신뢰구간, $z$/$\chi^2$/$t$) & 40 \\
2:00 & \textit{휴식} & & 10 \\
2:10 & 기초 통계 3: 상관과 회귀 & \Wc{6–8} (회귀 효과, RMSE, 잔차도) & 50 \\
3:00 & 검정 고르는 법 & \Wc{12} 검정 도구상자 (표 + 관심 있는 카드만) & 30 \\
3:30 & 응용 1: 통제와 interaction & \Hc{3, 6} + 각 장 끝 ``pandas 코드'' 직접 실행 & 50 \\
4:20 & \textit{휴식} & & 10 \\
4:30 & 응용 2: 재확인 & \Hc{7} logistic, \Hc{8} Bayesian (8.4 ``왜''는 꼭) & 40 \\
5:10 & 결론과 한계 & \Hc{9}, \Hc{11} 체크리스트, \Ac{표지 결론표} & 20 \\
5:30 & 도구 실행과 발표 문장 연습 & \code{algo\_toolkit.py} 실행, \Hc{9.2} 문장을 내 말로 & 30 \\
(다음 날) & 코드로 손에 익히기 & \Pc{0–12} 빈칸 노트북 (\code{stats\_basics\_practice.ipynb}) & 120 \\
\bottomrule
\end{tabular}
\end{center}
\aside{※ 범위 밖 (학부 1–2학년 목표를 넘는 것): sklearn을 쓰는 \Ac{6–7}의 결정트리와 \code{advanced/} 폴더의 분류 워크북은 나중에 선택으로.}\\
\aside{※ 건너뛰어도 되는 곳 (나중에 참고용): \Wc{9–11, 13}, \Hc{4, 10}, \Ac{1–10} 본문, 모든 연습 문제. 시간이 남으면 \Wc{15} 연습 문제 1–8번부터.}\\
\aside{※ 각 장은 새 페이지에서 시작하고, 목차의 항목을 누르면 그 쪽으로 이동합니다.}
 으로 넣는다.
\providecommand{\code}[1]{\texttt{#1}}
\newcolumntype{L}[1]{>{\raggedright\arraybackslash}p{#1}}
% 책별 장 표시 (작은 색 주석)
\newcommand{\Wc}[1]{{\scriptsize\textsf{\color{blue!60!black}[흐름 #1]}}}
\newcommand{\Hc}[1]{{\scriptsize\textsf{\color{orange!70!black}[응용 #1]}}}
\newcommand{\Ac}[1]{{\scriptsize\textsf{\color{green!45!black}[알고 #1]}}}
\newcommand{\Pc}[1]{{\scriptsize\textsf{\color{red!60!black}[실습 #1]}}}

\section*{먼저 읽기: 분석 전체 흐름과 책 안내}
\addcontentsline{toc}{section}{먼저 읽기: 분석 전체 흐름과 책 안내}

\textbf{네 권의 역할}
\begin{center}\small
\begin{tabular}{@{}l L{5.4cm} L{7.4cm}@{}}
\toprule
표시 & 책 (파일) & 역할 \\
\midrule
\Wc{} & pandas 흐름으로 따라가는 통계 (\code{pandas-stats-walkthrough}) & \textbf{기초}: 평균·SD·정규·상관·회귀·검정을 pandas 순서대로 \\
\Hc{} & 채용 데이터로 다시 배우는 통계 (\code{hiring-statistics-book}) & \textbf{응용}: 통제 회귀, logistic, Bayesian, 공정성으로 결론 내기 \\
\Ac{} & 알고리즘으로 보는 pandas (\code{algorithms-for-datathon}) & \textbf{도구}: 빠르고 정확한 코드, 결정트리·greedy 보조 분석 \\
\Pc{} & numpy·pandas로 직접 만드는 통계 (\code{stats-basics-workbook}) & \textbf{실습}: 빈칸 노트북으로 통계 공식을 직접 구현, 필요한 곳에 알고리즘 적용 \\
\bottomrule
\end{tabular}
\end{center}
\aside{※ 아래 표의 작은 색 주석이 ``이 단계는 어느 책 몇 장''입니다. 예: \Wc{5} = 흐름 책 5장, \Hc{6.9} = 응용 책 6.9절, \Ac{4} = 알고리즘 책 4장. 0장 = 응용 책 앞의 ``데이터 미리보기와 pandas 작동 원리''.}

\vspace{0.4em}
\textbf{Datathon 분석 흐름: 데이터를 받은 순간부터 발표까지}
\begin{center}\small
\renewcommand{\arraystretch}{1.25}
\begin{tabular}{@{}c L{3.3cm} L{4.6cm} L{5.6cm}@{}}
\toprule
단계 & 하는 일 & 대표 pandas 코드 & 해당 장 \\
\midrule
① & 질문을 통계 질문으로 & -- &
  \Hc{1} \Wc{1 추정·가설·검증 상자} \\
② & 불러오고 구조 확인 & \code{read\_csv}, \code{head}, \code{info} &
  \Hc{0.1–0.3} \Wc{1} \\
③ & 변수 분류 (결과 $Y$ / 보호속성 $A$ / 자격 $X$) & \code{dtypes}, \code{value\_counts} &
  \Hc{2} \Wc{1.1} \\
④ & 결측·이탈값 처리 & \code{isna}, \code{dropna}, \code{quantile} &
  \Hc{0.7, 2.4} \Wc{1.3, 2.3} \\
⑤ & 분포 요약 (중심, 퍼짐, 정규성) & \code{describe}, \code{std}, \code{(x-x.mean())/x.std()} &
  \Wc{2, 3} \Hc{3.1} \Ac{3 분위수} \Pc{1–3} \\
⑥ & 그룹별 합격률과 신뢰구간 & \code{groupby().mean()} &
  \Hc{0.5, 3.2} \Wc{4} \Ac{4 groupby 내부} \Pc{4} \\
⑦ & 교란 변수 확인 (Simpson) & \code{crosstab}, \code{groupby([a,b])} &
  \Hc{3.3} \Wc{12 ⑧} \\
⑧ & 가설검정 (차이가 우연인가) & \code{proportions\_ztest}, \code{chi2\_contingency}, \code{ttest\_ind} &
  \Wc{5, 11, 12} \Hc{5} \Pc{5–7} \\
⑨ & 상관과 단순회귀 & \code{corr}, \code{pd.cut}, \code{smf.ols} &
  \Wc{6–8} \Ac{5 pd.cut} \Pc{8–9} \\
⑩ & 통제 회귀와 interaction (핵심) & \code{smf.ols("y \~{} A + X")}, \code{A*B} &
  \Hc{0.8–0.9, 6} \Wc{13} \Pc{10} \\
⑪ & 다른 방법으로 재확인 & \code{smf.logit}, \code{rng.beta}, \code{DecisionTreeClassifier} &
  \Hc{7, 8} \Wc{9 bootstrap} \Pc{6, 11} \\
⑫ & 공정성 지표와 한계 & \code{rate["F"]/rate["M"]}, \code{idxmin} &
  \Hc{9} \Ac{8 minimax} \\
⑬ & 정리와 발표 & -- &
  \Hc{9.2 결론 문장, 11 체크리스트} \Wc{14 요약} \Ac{11} \\
\bottomrule
\end{tabular}
\end{center}
\aside{※ 이론이 막히면: 추정량·MoM·MLE \Wc{9–10} \Hc{4}, 가설 세우기·검정력 \Wc{11}, 코드가 느리면 \Ac{1–2, 4}. 한 번에 돌려 보기: \code{python3 algo\_toolkit.py hiring\_practice.csv} (⑥⑦⑨⑪⑫ 요약).}

\clearpage
\textbf{6시간 읽기 계획 (읽기 편한 순서)}
\begin{center}\small
\renewcommand{\arraystretch}{1.25}
\begin{tabular}{@{}l L{4.6cm} L{6.6cm} r@{}}
\toprule
시간 & 무엇을 & 어디를 & 분 \\
\midrule
0:00 & 데이터와 pandas 감 잡기 & \Hc{0} 전체, \Hc{1} 전체 지도 & 30 \\
0:30 & 기초 통계 1: 요약과 분포 & \Wc{1–3} (\code{describe}, 표준화, 정규근사) & 50 \\
1:20 & 기초 통계 2: 추정과 검정 & \Wc{4–5} (신뢰구간, $z$/$\chi^2$/$t$) & 40 \\
2:00 & \textit{휴식} & & 10 \\
2:10 & 기초 통계 3: 상관과 회귀 & \Wc{6–8} (회귀 효과, RMSE, 잔차도) & 50 \\
3:00 & 검정 고르는 법 & \Wc{12} 검정 도구상자 (표 + 관심 있는 카드만) & 30 \\
3:30 & 응용 1: 통제와 interaction & \Hc{3, 6} + 각 장 끝 ``pandas 코드'' 직접 실행 & 50 \\
4:20 & \textit{휴식} & & 10 \\
4:30 & 응용 2: 재확인 & \Hc{7} logistic, \Hc{8} Bayesian (8.4 ``왜''는 꼭) & 40 \\
5:10 & 결론과 한계 & \Hc{9}, \Hc{11} 체크리스트, \Ac{표지 결론표} & 20 \\
5:30 & 도구 실행과 발표 문장 연습 & \code{algo\_toolkit.py} 실행, \Hc{9.2} 문장을 내 말로 & 30 \\
(다음 날) & 코드로 손에 익히기 & \Pc{0–12} 빈칸 노트북 (\code{stats\_basics\_practice.ipynb}) & 120 \\
\bottomrule
\end{tabular}
\end{center}
\aside{※ 범위 밖 (학부 1–2학년 목표를 넘는 것): sklearn을 쓰는 \Ac{6–7}의 결정트리와 \code{advanced/} 폴더의 분류 워크북은 나중에 선택으로.}\\
\aside{※ 건너뛰어도 되는 곳 (나중에 참고용): \Wc{9–11, 13}, \Hc{4, 10}, \Ac{1–10} 본문, 모든 연습 문제. 시간이 남으면 \Wc{15} 연습 문제 1–8번부터.}\\
\aside{※ 각 장은 새 페이지에서 시작하고, 목차의 항목을 누르면 그 쪽으로 이동합니다.}

\clearpage
\renewcommand{\contentsname}{목차 (Contents)}
\tableofcontents

%=================================================
\section{Big-O와 RAM 모델: 왜 반복문을 쓰지 말라고 하나}
\aref{04 Big-O, 06 RAM model, 03 Why constants matter less}

\subsection{같은 $\Theta(n)$인데 230배 차이}
\begin{codebox}
x = np.random.normal(size=1\_000\_000)\\
\# (a) 파이썬 반복문 \hfill 0.0328초\\
t = 0.0\\
for v in x: t += v\\
\# (b) 벡터화 \hfill 0.00014초\\
x.mean()
\end{codebox}
두 방법 모두 원소를 한 번씩 보는 $\Theta(n)$입니다. 그런데 \textbf{약 230배} 차이가 납니다.

\begin{keybox}{강의의 ``상수는 덜 중요하다''와 충돌하나?}
\begin{itemize}
\item 강의: $n$이 커지면 $\Theta(n)$과 $\Theta(n^2)$의 차이가 상수 배 차이를 압도한다. \textbf{여전히 맞습니다} (아래 4장의 1400배).
\item 하지만 \emph{같은 차수끼리}는 상수가 전부입니다. RAM 모델은 ``기본 연산 1회 = 상수 시간''을 가정하는데, 파이썬에서는 \code{t += v} 한 번이 타입 확인, 객체 생성, 참조 카운트 등 수십 개의 기계 명령입니다. numpy는 같은 루프를 C로, 연속 메모리 위에서, SIMD로 돌립니다.
\end{itemize}
\aside{※ 그래서 pandas의 제1원칙: \textbf{행 단위 반복을 열 단위 연산으로 바꿔라 (vectorization).}}
\end{keybox}

\subsection{pandas에서의 실제 예}
\begin{center}
\begin{tabular}{@{}l r r@{}}
\toprule
\code{test\_score}의 z-score (582행) & 시간 & 배율 \\
\midrule
\code{for \_, r in d.iterrows(): ...} & 5.1 ms & 110배 \\
\code{(d.test\_score - d.test\_score.mean()) / d.test\_score.std()} & 0.046 ms & 1 \\
\bottomrule
\end{tabular}
\end{center}
\aside{※ \code{iterrows}는 행마다 Series 객체를 새로 만듭니다. \code{apply(axis=1)}도 비슷하게 느립니다. 600행에서는 체감이 없지만, 수십만 행 datathon 데이터에서는 몇 초와 몇 분의 차이가 됩니다.}

\begin{dtbox}
\begin{itemize}
\item 열 연산 (\code{+ - * /}, \code{np.where}, \code{.str}, \code{.dt}), \code{groupby}, \code{merge}를 우선 쓰고, 반복문은 마지막 수단으로.
\item 코드가 느리면 먼저 ``이 루프의 차수가 $n^2$인가?''(4장)를 확인하고, 그다음 ``파이썬 루프인가?''를 확인합니다.
\end{itemize}
\end{dtbox}

%=================================================
\section{정렬: \code{sort\_values} 안에서 일어나는 일}
\aref{10 Merge sort, 13 Heap sort, 14–15 Quick sort, 16 Lower bound, 17 Counting/Radix}

\subsection{\code{kind=} 옵션은 강의 목차 그 자체}
\begin{codebox}
df.sort\_values("test\_score", kind="quicksort") \hfill \# 기본값\\
\# kind: \{"quicksort", "mergesort", "heapsort", "stable"\}
\end{codebox}
\begin{center}
\begin{tabular}{@{}l p{8.6cm} l c@{}}
\toprule
\code{kind} & 실제 알고리즘 (numpy 문서) & 최악 & stable \\
\midrule
\code{quicksort} & \textbf{introsort}: quick sort, 재귀가 깊어지면 heap sort로 전환 & $O(n\log n)$ & \ding{55} \\
\code{heapsort} & heap sort & $O(n\log n)$ & \ding{55} \\
\code{mergesort}, \code{stable} & 실수: \textbf{timsort} (merge + insertion), 작은 정수: \textbf{radix sort} & $O(n\log n)$ & \ding{51} \\
\bottomrule
\end{tabular}
\end{center}
\aside{※ Introsort = 강의 14–15의 quick sort (평균 $\Theta(n\log n)$, 최악 $\Theta(n^2)$) + 강의 13의 heap sort (최악 $\Theta(n\log n)$). Quick sort가 나쁜 pivot을 계속 고르면 heap sort로 바꿔서, \emph{평균은 quick sort만큼 빠르고 최악은 heap sort만큼 안전}합니다. 두 강의를 합친 실전형입니다.}

\subsection{측정: 100만 개}
\begin{center}
\begin{tabular}{@{}l l r@{}}
\toprule
데이터 & 방법 & 시간 \\
\midrule
무작위 실수 & quicksort (introsort) & 21 ms \\
 & mergesort / stable (timsort) & 68 / 70 ms \\
 & heapsort & 74 ms \\
\midrule
이미 정렬된 실수 & quicksort & 20 ms \\
 & stable (timsort) & \textbf{0.9 ms} \\
\midrule
0–99 정수 (\code{int8}) & quicksort & 19 ms \\
 & stable (radix sort) & \textbf{0.7 ms} \\
\bottomrule
\end{tabular}
\end{center}

\begin{keybox}{표에서 배우는 것}
\begin{enumerate}
\item \textbf{무작위 자료에서는 quicksort가 가장 빠름.} 셋 다 $\Theta(n\log n)$이지만 quick sort는 메모리를 순서대로 읽어서 (cache-friendly) 상수가 작습니다. 그래서 기본값입니다.
\item \textbf{Heap sort가 가장 느림.} 최악 보장은 좋지만, 힙에서 부모–자식 사이를 건너뛰며 접근해서 캐시 효율이 나쁩니다. 그래서 단독보다는 introsort의 안전장치로 쓰입니다.
\item \textbf{정렬된 자료에서 timsort가 23배 빠름.} 이미 정렬된 구간(run)을 찾아 합치므로 거의 $\Theta(n)$. Insertion sort의 best case $\Theta(n)$ \aref{9}과 같은 원리입니다. 시계열처럼 ``거의 정렬된'' 자료에 유리합니다.
\item \textbf{작은 정수에서 radix sort가 27배 빠름.} 비교 정렬의 하한 $\Omega(n\log n)$ \aref{16}은 \emph{비교만 쓰는} 정렬에 대한 것입니다. Counting/radix sort는 값을 직접 인덱스로 쓰므로 $\Theta(n+k)$, $\Theta(d(n+k))$로 그 하한을 피합니다 \aref{17}.
\end{enumerate}
\end{keybox}

\subsection{Stability가 실제로 필요한 순간}
\term{Stable sort}: 키가 같은 원소들의 원래 순서를 유지합니다. \textbf{여러 키로 정렬할 때} 중요합니다.
\begin{codebox}
\# 부서별로, 같은 부서 안에서는 점수 내림차순\\
d.sort\_values("test\_score", ascending=False)\textbackslash\\
\ \ .sort\_values("department", kind="stable") \hfill \# 2단계 정렬: 두 번째는 반드시 stable\\
d.sort\_values(["department", "test\_score"], ascending=[True, False]) \hfill \# 같은 결과, 권장
\end{codebox}
\aside{※ 두 번째 정렬이 unstable이면 같은 부서 안의 점수 순서가 깨질 수 있습니다. Radix sort가 \emph{낮은 자리부터} stable하게 정렬하는 원리 \aref{17}와 정확히 같습니다: 덜 중요한 키 먼저, 더 중요한 키를 나중에 stable하게.}

\subsection{통계에서 정렬이 쓰이는 곳}
\begin{itemize}
\item \term{Order statistics} $x_{(1)}\le\dots\le x_{(n)}$ = 정렬 결과 그 자체 \sref{2장} (강의노트 Def 2.1).
\item \term{Empirical CDF}, \term{Q-Q plot}: 정렬된 자료를 이론 분위수와 짝지음 \sref{3장}.
\item \term{Rank}: \code{s.rank()} = 정렬 후 위치. Mann–Whitney, Spearman 상관 같은 비모수 검정의 재료.
\item \code{drop\_duplicates}, \code{merge\_asof}, 시계열 리샘플링은 정렬된 상태를 가정하거나 내부에서 정렬합니다.
\end{itemize}

%=================================================
\section{선택(selection)과 힙: median, quantile, top-$k$}
\aref{13 Heap sort, 14–15 Quick sort (partition)}

\subsection{중앙값을 구하려고 전부 정렬할 필요는 없다}
$k$번째 작은 값만 필요하면 \term{quickselect}를 씁니다. Quick sort의 \code{Partition}을 한 번 하고, $k$가 있는 \emph{한쪽만} 재귀합니다.
\begin{align*}
\text{Quick sort:}\quad & T(n)=2T(n/2)+\Theta(n)=\Theta(n\log n),\\
\text{Quickselect:}\quad & T(n)=T(n/2)+\Theta(n)=\Theta(n)\ (\text{평균}).
\end{align*}
\aside{※ 두 번째 점화식은 $n+n/2+n/4+\dots\le2n$. 강의 18–19의 recursion tree로 바로 풀립니다.}
\begin{center}
\begin{tabular}{@{}l r@{}}
\toprule
100만 개의 중앙값 & 시간 \\
\midrule
\code{np.sort(x)[n//2]} (전부 정렬) & 21 ms \\
\code{np.partition(x, n//2)[n//2]} (introselect) & 7.3 ms \\
\code{np.median(x)} (내부에서 partition 사용) & 9.0 ms \\
\bottomrule
\end{tabular}
\end{center}
\aside{※ \code{np.quantile}, \code{Series.quantile}, \code{describe()}의 25/50/75\%도 내부에서 \code{partition}을 씁니다 (소스에서 확인). 즉 통계책 2장의 사분위수 계산이 quickselect입니다.}

\subsection{Top-$k$: \code{nlargest}}
\begin{center}
\begin{tabular}{@{}l l r@{}}
\toprule
100만 개 중 상위 10개 & 방법 & 시간 \\
\midrule
\code{s.sort\_values(ascending=False).head(10)} & 전부 정렬 $\Theta(n\log n)$ & 77 ms \\
\code{heapq.nlargest(10, x)} & 크기 $k$ 힙 $\Theta(n\log k)$ (파이썬) & 24 ms \\
\code{s.nlargest(10)} & selection + 10개만 정렬 & \textbf{10 ms} \\
\bottomrule
\end{tabular}
\end{center}
\begin{itemize}
\item \textbf{힙 방법}: 크기 $k$인 min-heap을 유지. 새 원소가 힙의 최솟값(루트)보다 크면 교체하고 \code{HeapifyDown} $O(\log k)$. 강의 13의 힙 연산 그대로입니다.
\item \textbf{pandas 방법} (소스 확인): \code{kth\_smallest}(quickselect 계열)로 $k$번째 값을 찾고, 그보다 큰 원소만 골라 merge sort.
\end{itemize}
\aside{※ 힙은 ``전체 정렬은 필요 없고 가장 큰/작은 것만 반복해서 꺼내고 싶을 때''의 자료구조입니다. 이 성질 때문에 6장 결정트리(best-first)와 9장 Dijkstra에서 priority queue로 다시 나옵니다.}

\begin{dtbox}
\begin{itemize}
\item ``점수 상위 10명'', ``합격률 낮은 부서 3개''는 \code{nlargest}/\code{nsmallest}.
\item 분위수, 중앙값, IQR은 정렬 없이 빠르게 계산되므로 큰 데이터에서도 \code{describe()}를 부담 없이 씁니다.
\end{itemize}
\end{dtbox}

%=================================================
\section{해싱과 counting sort: \code{groupby}와 \code{merge}}
\aref{17 Counting sort}

\subsection{\code{groupby} 내부}
\code{df.groupby("gender").hired.mean()}은 대략 세 단계입니다.
\begin{enumerate}
\item \textbf{Factorize (해싱)}: \code{"F","M"} $\to$ 정수 코드 $0,1$. 해시 테이블로 평균 $O(n)$.
\item \textbf{Counting sort}: 코드별로 행을 모음. pandas 소스의 주석 그대로 ``\code{groupsort\_indexer} implements counting sort''. 그룹 수 $k$가 작으므로 $\Theta(n+k)$.
\item 그룹별로 합과 개수를 누적해 평균 계산.
\end{enumerate}
\aside{※ 강의 17의 counting sort: (1) 각 키의 개수를 세고 (2) 누적합으로 시작 위치를 정하고 (3) 원소를 제자리에 놓는다. \code{groupby}가 정확히 이렇게 그룹의 시작/끝 위치를 만듭니다. 그래서 600행 groupby가 0.09 ms.}

\subsection{\code{merge}: $O(nm)$ vs $O(n+m)$}
\begin{center}
\begin{tabular}{@{}l l r@{}}
\toprule
3000행 $\times$ 3000행 조인 & 복잡도 & 시간 \\
\midrule
이중 for문 (모든 쌍 비교) & $\Theta(nm)$ & 448 ms \\
\code{a.merge(b, on="k")} (hash join) & $\Theta(n+m)$ 평균 & 0.32 ms \\
\bottomrule
\end{tabular}
\end{center}
\textbf{1400배.} 1장과 달리 이번에는 \emph{차수}의 차이입니다. 행이 10배 늘면 이중 루프는 100배, hash join은 10배 느려집니다.\\
\aside{※ Hash join: 작은 쪽 키로 해시 테이블을 만들고 ($O(m)$), 다른 쪽을 훑으며 조회 ($O(n)$, 조회당 평균 $O(1)$).}

\begin{dtbox}
여러 테이블(지원자, 면접, 부서 정보)을 받으면 반드시 \code{merge}로 합칩니다. 합친 뒤 \code{len()}이 예상과 같은지 확인하세요. 키가 중복되면 행이 곱으로 늘어납니다 (\code{validate="one\_to\_one"} 옵션).
\end{dtbox}

%=================================================
\section{이진 탐색: \code{searchsorted}와 \code{pd.cut}}
\aref{07 Search, Case study 1.1 Binary search}

\subsection{정렬되어 있으면 $O(\log n)$}
\begin{center}
\begin{tabular}{@{}l l r@{}}
\toprule
정렬된 100만 개에서 1000개 값의 위치 찾기 & 복잡도 & 시간 \\
\midrule
선형 탐색 (\code{np.argmax(xs >= v)}) & $\Theta(n)$ 각각 & 79 ms \\
\code{np.searchsorted(xs, q)} & $\Theta(\log n)$ 각각 & 0.19 ms \\
\bottomrule
\end{tabular}
\end{center}
\aside{※ 강의 case study의 전제조건 ``배열 $A$가 정렬되어 있음''이 그대로 적용됩니다. \code{searchsorted}에 정렬되지 않은 배열을 넣으면 \emph{오류 없이 틀린 답}이 나옵니다. 정확성 증명의 전제조건이 실전에서 왜 중요한지 보여 주는 예입니다.}

\subsection{어디에 숨어 있나}
\begin{itemize}
\item \code{pd.cut(score, bins)}: 각 점수가 어느 구간인지 = 경계값 배열에서 이진 탐색 (소스에서 \code{searchsorted} 확인). 통계책 7장 ``평균의 그래프''에서 쓴 바로 그 함수입니다 \sref{7장}.
\item \code{pd.merge\_asof}: ``가장 가까운 이전 시점''과 조인. 양쪽이 정렬되어 있어야 함.
\item 정렬된 인덱스에서 \code{df.loc["2024-01":"2024-03"]} 같은 범위 선택.
\item \code{np.quantile}의 역함수 격인 ``이 점수는 몇 백분위?'' $=$ \code{np.searchsorted(np.sort(s), 85) / len(s)} (통계책 3장의 95.7\%).
\end{itemize}

%=================================================
\section{트리, DFS, BFS: 결정트리와 연결요소}
\aref{12 Trees, 13 Heap (완전이진트리)}

\subsection{결정트리 = 데이터로 만든 이진트리}
\begin{codebox}
from sklearn.tree import DecisionTreeClassifier, export\_text\\
X = pd.get\_dummies(d[["gender","department","degree","test\_score","years\_experience"]],\\
\ \ \ \ \ \ \ \ \ \ \ \ \ \ \ \ \ \ \ drop\_first=True).astype(float)\\
clf = DecisionTreeClassifier(max\_depth=3, random\_state=0).fit(X, d.hired)\\
print(export\_text(clf, feature\_names=list(X.columns)))
\end{codebox}
\begin{codebox}
|--- department\_Marketing <= 0.50 \hfill (Engineering)\\
|\ \ \ |--- test\_score <= 68.50\\
|\ \ \ |\ \ \ |--- years\_experience <= 6.50 \ \ -> class 0\\
|\ \ \ |\ \ \ |--- years\_experience > 6.50 \ \ \ -> class 1\\
|\ \ \ |--- test\_score > 68.50\\
|\ \ \ |\ \ \ |--- gender\_M <= 0.50 \ \ \ \ \ \ \ \ \ -> class 0 \hfill (여성)\\
|\ \ \ |\ \ \ |--- gender\_M > 0.50 \ \ \ \ \ \ \ \ \ \ -> class 1 \hfill (남성)\\
|--- department\_Marketing > 0.50 \hfill (Marketing)\\
|\ \ \ |--- test\_score <= 79.50 ... -> class 0
\end{codebox}
\textbf{해석}: 트리가 스스로 ``Engineering + 점수 68.5 초과'' 그룹에서 \emph{성별}로 분할했습니다. 점수가 높은 Engineering 지원자 중 남성은 합격, 여성은 불합격으로 예측됩니다. 통계책·이전 책에서 회귀 interaction으로 찾은 ``Engineering에서의 성별 격차''를 전혀 다른 방법이 다시 발견한 것입니다.\\
\aside{※ Feature importance: Marketing 0.44, 점수 0.34, 경력 0.18, 성별 0.05, 학위 0. 다만 트리의 importance는 효과의 \emph{크기나 유의성}이 아니므로 결론은 회귀로 내리고, 트리는 ``패턴 탐색''용으로 씁니다.}

\subsection{트리를 만들고 읽는 순서 = DFS / best-first}
sklearn 소스(\code{\_tree.pyx})에는 두 가지 builder가 있습니다.
\begin{center}
\begin{tabular}{@{}l p{3.4cm} p{5.6cm}@{}}
\toprule
builder & 언제 & 자료구조 = 탐색 방식 \\
\midrule
\code{DepthFirstTreeBuilder} & 기본 (\code{max\_depth} 등) & \textbf{스택} = DFS \\
\code{BestFirstTreeBuilder} & \code{max\_leaf\_nodes} 지정 시 & \textbf{우선순위 큐 (힙)} = 불순도 감소가 가장 큰 노드부터 \\
\bottomrule
\end{tabular}
\end{center}
\begin{itemize}
\item \code{export\_text}의 출력 순서는 \textbf{전위 순회 (preorder DFS)}: 노드 $\to$ 왼쪽 서브트리 전체 $\to$ 오른쪽 서브트리.
\item \code{max\_leaf\_nodes=4}로 바꾸면 힙에서 ``가장 이득이 큰 잎''부터 꺼내 분할합니다. 결과: Marketing $\to$ 0; Engineering은 점수 68.5로 나누고, 그 아래 낮은 점수 쪽만 경력으로 한 번 더 분할. 잎이 4개뿐이라 성별 분할은 우선순위에서 밀립니다.
\item 노드 수가 $n$인 균형 트리의 깊이는 $\Theta(\log n)$ \aref{12}. \code{max\_depth=3}이면 잎이 최대 $2^3=8$개.
\end{itemize}
\aside{※ BFS (level-order)는 같은 깊이의 노드를 먼저 모두 방문합니다. 트리를 그림으로 그릴 때 층별 배치가 BFS 순서입니다.}

\subsection{BFS/DFS가 직접 필요한 경우: 연결요소}
데이터가 \emph{관계}로 주어지면 그래프 탐색이 필요합니다. 예: 같은 이메일이나 전화번호를 공유하는 지원 기록을 한 사람으로 묶기 (중복 지원 제거).
\begin{codebox}
from scipy.sparse.csgraph import connected\_components\\
\# A: 기록 i와 j가 같은 이메일/전화를 공유하면 1인 인접행렬 (희소)\\
n\_people, label = connected\_components(A, directed=False)\\
df["person\_id"] = label
\end{codebox}
\aside{※ \code{connected\_components}는 내부적으로 BFS/DFS를 반복합니다: 방문하지 않은 노드에서 탐색을 시작해 닿는 모든 노드에 같은 라벨. $\Theta(V+E)$. 직접 짜면 \code{collections.deque}로 BFS를 구현합니다.}\\
\aside{※ 중첩 JSON을 표로 펴는 \code{pd.json\_normalize}도 재귀적 DFS입니다.}

\begin{dtbox}
이번 데이터(한 행 = 한 지원자, 관계 없음)에서는 그래프 탐색이 필요 없습니다. 다만 결정트리를 \emph{보조 분석}으로 써서 ``회귀가 놓친 interaction이 있는지'' 찾는 데는 유용합니다. 여기서도 Engineering $\times$ 성별을 바로 짚어냈습니다.
\end{dtbox}

%=================================================
\section{Heuristic과 greedy: 최적해 대신 좋은 해}

\subsection{결정트리는 greedy heuristic}
각 노드에서 ``지금 불순도(Gini)를 가장 많이 줄이는 분할 하나''만 고르고 되돌아가지 않습니다. 전체적으로 최적인 트리를 찾는 문제는 NP-hard라서, greedy가 현실적인 대안입니다.
\begin{itemize}
\item 각 노드에서 변수마다 값을 \textbf{정렬}해 가능한 분할점을 모두 훑습니다. sklearn 소스(\code{\_partitioner.pyx})에서 쓰는 정렬은 \textbf{introsort} (2장). 그래서 트리 학습 비용은 대략 $O(d\cdot n\log n)$ per level입니다.
\end{itemize}

\subsection{변수 선택: forward stepwise (greedy)}
변수 조합은 $2^p$개라 다 볼 수 없습니다. AIC (통계책 이전 책 §7)를 \textbf{평가 함수}로 놓고, 매 단계 AIC를 가장 많이 줄이는 변수 하나를 추가합니다.
\begin{center}\small
\begin{tabular}{@{}c l l l@{}}
\toprule
단계 & 후보별 AIC (logistic) & 선택 & AIC \\
\midrule
0 & (절편만) & -- & 735.8 \\
1 & 부서 679.4, 점수 701.2, 성별 706.5, 경력 715.1, 학위 731.2 & 부서 & 679.4 \\
2 & 점수 638.4, 경력 667.9, 성별 675.4, 학위 677.3 & 점수 & 638.4 \\
3 & 경력 625.9, 성별 635.8, 학위 635.9 & 경력 & 625.9 \\
4 & 학위 623.3, 성별 623.5 & 학위 & 623.3 \\
5 & 성별 620.7 & 성별 & 620.7 \\
\bottomrule
\end{tabular}
\end{center}
\textbf{해석}: 성별은 마지막에 들어오지만 그래도 AIC를 낮춥니다 (623.3 $\to$ 620.7). 다른 모든 자격 변수를 넣은 뒤에도 성별이 설명력을 더한다는 뜻으로, 통계책 13장의 ``통제 후 성별 계수 유의''와 같은 결론입니다.\\
\aside{※ Greedy의 한계: 단계 1에서 성별(706.5)은 점수(701.2)와 비슷했지만 부서와 겹치는 정보라 뒤로 밀렸습니다. 순서가 결과 해석을 좌우하므로, stepwise는 탐색용이고 최종 모형은 이론(무엇을 통제해야 하는가)으로 정합니다.}

\subsection{Heuristic function과의 관계}
\begin{itemize}
\item A* 탐색의 heuristic function $h(n)$: 목표까지 남은 비용의 \emph{추정치}. 정확하지 않아도 탐색 방향을 잡아 줍니다. Admissible ($h\le$ 실제 비용)이면 최적해 보장.
\item 데이터 분석의 ``heuristic'': Gini 감소, AIC, validation 오차 같은 \emph{평가 함수}로 거대한 탐색 공간(트리 구조, 변수 조합, 하이퍼파라미터)을 greedy하게 훑는 것.
\item 하이퍼파라미터 탐색: \code{GridSearchCV} (전수 탐색), \code{RandomizedSearchCV} (무작위), Bayesian optimization (모형 기반 heuristic).
\end{itemize}

%=================================================
\section{Minimax: 직접은 없지만 같은 생각이 있다}

\subsection{원래 의미}
게임 트리에서 내 차례에는 max, 상대 차례에는 min을 취해 \emph{최악의 상대를 가정했을 때 가장 좋은 수}를 고릅니다. 구현은 DFS + 깊이 제한 + heuristic 평가 함수 (+ alpha-beta 가지치기). 데이터 분석에 게임 트리가 나올 일은 거의 없습니다.

\subsection{같은 원리가 나오는 곳}
\textbf{1. Minimax 추정량 (통계적 결정이론).} 추정량의 위험(MSE)은 진짜 $\theta$에 따라 달라집니다 (통계책 9장). ``$\theta$가 최악일 때의 MSE를 가장 작게'' 만드는 추정량이 minimax 추정량입니다:
\[
\hat\theta_{\text{minimax}}=\arg\min_{\hat\theta}\ \max_{\theta}\ \text{MSE}(\hat\theta;\theta).
\]
\aside{※ 자연(nature)을 ``나를 괴롭히는 상대''로 보는 게임입니다. Bayesian은 반대로 $\theta$에 대한 prior로 \emph{평균} 위험을 최소화합니다.}

\textbf{2. Minimax 공정성 (fairness).} ``가장 불리한 그룹의 성과를 최대화''하는 기준입니다:
\[
\max_{\text{정책}}\ \min_{\text{그룹 }g}\ \text{성과}_g.
\]
\begin{codebox}
g = df.groupby(["department","gender"]).hired.mean()\\
g.idxmin(), g.min() \hfill \# (Marketing, M), 0.098\\
g.idxmax(), g.max() \hfill \# (Engineering, M), 0.527
\end{codebox}
\aside{※ 평균 격차(demographic parity)만 보면 놓치는 것: 가장 불리한 셀은 여성이 아니라 \emph{Marketing 남성} (9.8\%)입니다. 다만 61명뿐이라 불확실성이 큽니다 (통계책 11장 검정력). 발표에서 ``최악 그룹''을 함께 보여 주면 분석의 깊이가 드러납니다.}

\textbf{3. Robust 결론.} 여러 모형(LPM, logistic, Bayesian, 결정트리) 중 \emph{가장 불리한 모형에서도} 결론이 유지되는가를 보는 것도 minimax식 사고입니다. 이번 데이터에서는 Engineering 성별 격차가 모든 방법에서 나왔습니다.

%=================================================
\section{Dijkstra: 이번 datathon에는 필요 없다}

\subsection{언제 쓰나}
가중치가 음수가 아닌 그래프에서 한 출발점으로부터 모든 노드까지의 최단 경로. 데이터 형태가 \textbf{노드와 간선}일 때만 필요합니다.
\begin{center}
\begin{tabular}{@{}l l@{}}
\toprule
datathon 주제 예 & 쓰임 \\
\midrule
교통·물류 (버스 노선, 배송) & 지역 간 최단 이동 시간, 접근성 지표 \\
소셜 네트워크 & 사람 사이 거리, 중심성 (closeness) \\
인프라 (전력망, 상수도) & 장애 시 우회 경로 \\
\bottomrule
\end{tabular}
\end{center}
\begin{codebox}
from scipy.sparse.csgraph import dijkstra \hfill \# 또는 networkx.shortest\_path\_length\\
dist = dijkstra(adjacency\_matrix, indices=source)
\end{codebox}

\subsection{알고리즘의 연결고리}
\begin{itemize}
\item 핵심 자료구조: \textbf{min-heap 우선순위 큐}. 매번 ``지금까지 거리가 가장 짧은 미방문 노드''를 꺼냅니다 (\code{ExtractMin}, $O(\log V)$). 강의 13의 힙 연산 그대로이고, 3장의 top-$k$, 6장의 best-first tree builder와 같은 도구입니다.
\item 복잡도: 이진 힙으로 $O((V+E)\log V)$.
\item 가중치가 모두 같으면 Dijkstra = \textbf{BFS} (6장). A*는 Dijkstra에 heuristic function $h$를 더한 것 (7장).
\end{itemize}
\aside{※ 이번 채용 데이터는 지원자 사이의 관계가 없는 표 데이터라 그래프가 만들어지지 않습니다. 그러니 이번 datathon에서는 넘겨도 됩니다.}

%=================================================
\section{분할정복과 점화식: 분석 파이프라인의 비용}
\aref{10 Merge sort, 11 Max subarray, 18–19 Recurrences, 02 Peak finding}

\subsection{Split–apply–combine}
\code{groupby().apply(f)}는 데이터를 그룹으로 \emph{나누고} (split), 각각 계산하고 (apply), 합칩니다 (combine). 분할정복의 divide–conquer–combine과 같은 구조입니다. 차이: 재귀가 아니라 한 단계라서 비용은 $\Theta(n)+\sum_g \text{cost}(f, n_g)$.

\subsection{재표본 방법의 비용}
\begin{center}
\begin{tabular}{@{}l l l@{}}
\toprule
방법 (통계책) & 비용 & 이번 데이터 \\
\midrule
Bootstrap $B$회 (9장) & $\Theta(B\cdot n)$ $\times$ 통계량 비용 & $2000\times600$, 수 초 \\
Permutation $P$회 (12장 ⑬) & $\Theta(P\cdot n)$ & $5000\times600$ \\
$k$-fold cross-validation & $k\times$ 모형 학습 비용 & -- \\
결정트리 학습 & 대략 $O(d\,n\log n\cdot\text{depth})$ & 순식간 \\
\bottomrule
\end{tabular}
\end{center}
\aside{※ $B$를 2배로 하면 시간도 2배, Monte Carlo 오차는 $1/\sqrt2$배. ``정밀도를 2배로 하려면 4배의 반복''이라는 $\sqrt{n}$ 법칙이 시뮬레이션 비용에도 적용됩니다.}

\subsection{강의의 다른 알고리즘이 쓰이는 곳}
\begin{itemize}
\item \textbf{Maximum subarray} \aref{11}: 시계열에서 ``누적 이득이 가장 큰 구간'' (예: 주가가 가장 많이 오른 기간, 합격률이 가장 높았던 연속 기간). Kadane 알고리즘 $\Theta(n)$. pandas로는 \code{cumsum}과 \code{cummin}의 조합으로 벡터화할 수 있습니다.
\item \textbf{Peak finding} \aref{02}: 시계열의 국소 최댓값. \code{scipy.signal.find\_peaks}.
\item \textbf{Binary search on answer}: ``합격률 차이가 0.8 비율(four-fifths)을 넘는 최소 점수 기준''처럼 단조인 조건의 경계 찾기.
\end{itemize}
\aside{※ 이번 채용 데이터에는 시간 축이 없어서 이 두 알고리즘은 쓰이지 않습니다.}

%=================================================
\section{정리: datathon 실전 규칙}
\begin{enumerate}
\item \textbf{행 반복 금지.} \code{iterrows}, \code{apply(axis=1)} 대신 열 연산 (1장, 110–230배).
\item \textbf{이중 루프 금지.} 두 표의 짝짓기는 \code{merge} (4장, 1400배, 차수의 차이).
\item \textbf{정렬은 기본값으로 충분.} 여러 키 정렬은 \code{sort\_values([k1, k2])} 한 번에. 단계별로 하면 \code{kind="stable"} (2장).
\item \textbf{상위 $k$개는 \code{nlargest}}, 분위수는 \code{quantile} (3장).
\item \textbf{구간 나누기는 \code{pd.cut}/\code{qcut}}, 정렬된 자료 탐색은 \code{searchsorted} (5장). 전제조건: 정렬.
\item \textbf{결정트리는 interaction 탐색용 보조 도구}, 결론은 회귀와 검정으로 (6장).
\item \textbf{Stepwise 선택은 탐색용}, 최종 모형은 이론으로 (7장).
\item \textbf{최악 그룹도 보고}: \code{groupby(...).mean().idxmin()} (8장).
\item 그래프 알고리즘 (BFS/DFS/Dijkstra)은 데이터가 \emph{관계}일 때만 (6, 9장).
\end{enumerate}

%=================================================
\section{연습 문제}
\begin{enumerate}
\item \code{df.sort\_values("test\_score")}의 기본 알고리즘은 무엇이고, 최악의 경우 시간복잡도는? 순수 quick sort와 무엇이 다른가?
\item 100만 개의 \code{years\_experience} (0–11 정수)를 정렬할 때 \code{kind="stable"}이 \code{quicksort}보다 훨씬 빠를 수 있는 이유를 정렬 하한과 연결해 설명하시오.
\item 부서 오름차순, 같은 부서에서는 점수 내림차순으로 정렬하려고 다음과 같이 썼다.
\begin{codebox}
d.sort\_values("department").sort\_values("test\_score", ascending=False, kind="stable")
\end{codebox}
 무엇이 틀렸는가?
\item 중앙값을 quickselect로 구하는 평균 시간이 $\Theta(n)$인 이유를 점화식으로 보이시오.
\item 크기 $k$인 min-heap으로 $n$개 중 상위 $k$개를 구하는 알고리즘을 쓰고, 복잡도가 $O(n\log k)$임을 보이시오. 왜 max-heap이 아니라 min-heap인가?
\item 지원자 표 (10만 행)와 면접 표 (10만 행)를 지원자 ID로 이중 for문으로 합치면 4장 측정(3000행에 0.45초)을 기준으로 대략 얼마나 걸리겠는가? \code{merge}로는?
\item 결정트리에서 \code{max\_depth=3}과 \code{max\_leaf\_nodes=4}는 각각 어떤 탐색 방식(자료구조)으로 트리를 키우는가?
\item Forward stepwise에서 성별이 마지막에 선택되었다. 이것이 ``성별은 가장 덜 중요하다''는 뜻인가?
\end{enumerate}

\section{풀이}
\textbf{1.}
\Ans Introsort, 최악 $O(n\log n)$. 순수 quick sort는 최악 $\Theta(n^2)$ (매번 최솟값/최댓값이 pivot)이지만, introsort는 재귀 깊이가 $c\log n$을 넘으면 heap sort로 전환해 최악을 막습니다.

\textbf{2.}
\Ans 비교 정렬의 하한 $\Omega(n\log n)$은 원소를 \emph{비교}로만 다룰 때의 하한입니다. 값의 범위가 작은 정수($k=12$)면 counting/radix sort로 개수를 세서 $\Theta(n+k)=\Theta(n)$에 정렬할 수 있습니다. numpy의 \code{stable}은 작은 정수형에 radix sort를 씁니다 (2장 측정: 27배).
\Just 단, \code{years\_experience}가 \code{int64}로 저장되어 있으면 radix 대신 timsort가 쓰일 수 있습니다. \code{astype("int8")}로 줄이면 이득을 봅니다.

\textbf{3.}
\Ans 순서가 거꾸로입니다. stable 다단계 정렬은 \emph{덜 중요한 키를 먼저}, 더 중요한 키를 \emph{나중에 stable하게} 정렬해야 합니다. 올바른 코드:
\begin{codebox}
d.sort\_values("test\_score", ascending=False).sort\_values("department", kind="stable")\\
d.sort\_values(["department", "test\_score"], ascending=[True, False]) \hfill \# 권장
\end{codebox}

\Just Radix sort가 낮은 자리부터 정렬하는 이유와 같습니다 \aref{17}.

\textbf{4.}
\Ans 평균적으로 partition이 배열을 비슷하게 나누면 한쪽만 재귀하므로 $T(n)=T(n/2)+cn$. 펼치면 $cn(1+\frac12+\frac14+\cdots)\le2cn=\Theta(n)$.
\Just Quick sort는 양쪽 모두 재귀해서 $T(n)=2T(n/2)+cn=\Theta(n\log n)$. 차이는 재귀 호출이 1개냐 2개냐입니다.

\textbf{5.}
\Ans 처음 $k$개로 min-heap을 만들고 ($O(k)$), 나머지 각 원소 $x$에 대해 $x>$ 루트이면 루트를 $x$로 바꾸고 HeapifyDown ($O(\log k)$). 총 $O(k+(n-k)\log k)=O(n\log k)$.
\Just 힙의 루트는 ``지금까지 상위 $k$개 중 가장 작은 값'', 즉 탈락 후보입니다. 새 원소와 비교할 대상이 이것이므로 min-heap이어야 루트에서 $O(1)$로 확인할 수 있습니다.

\textbf{6.}
\Ans 이중 루프는 $\Theta(nm)$이므로 $(10^5/3000)^2\approx1111$배 $\Rightarrow$ 약 $0.45\times1111\approx500$초 (8분 이상). \code{merge}는 $\Theta(n+m)$이므로 약 33배 $\Rightarrow$ 약 0.01초.

\textbf{7.}
\Ans \code{max\_depth}: \code{DepthFirstTreeBuilder}, 스택 기반 DFS. \code{max\_leaf\_nodes}: \code{BestFirstTreeBuilder}, 불순도 감소량을 키로 하는 우선순위 큐 (힙).

\textbf{8.}
\Ans 아닙니다. Greedy는 각 단계에서 \emph{그 시점의 추가 이득}만 봅니다. 성별의 정보 일부가 부서와 겹치므로 (성별 $\times$ 부서 $\chi^2=132$), 부서가 먼저 들어가면 성별의 추가 이득이 작아 보입니다. 그래도 마지막에 AIC를 낮췄다는 것은 \emph{모든 자격 변수를 통제한 후에도} 성별이 설명력을 가진다는 뜻이고, 이것이 bias 분석의 핵심 질문입니다.

\end{document}
